\documentclass[a4paper,12pt]{article}
\usepackage[utf8]{inputenc}
\usepackage[brazil]{babel}
\usepackage{enumitem}
\usepackage{hyperref}

\title{Produtos Oferecidos pela Azure, GCP, Alibaba Cloud e suas respectivas Funcionalidades}
\author{Breno Fernandes - 01169313}
\date{\today}

\begin{document}

\maketitle

\break

\section*{Azure}

\subsection*{Azure Virtual Machines (VMs)}
As VMs na Azure são usadas para hospedar servidores virtuais com configurações flexíveis, permitindo ao usuário escolher o sistema operacional, a CPU, a memória e o armazenamento.

\textbf{Tipos de VMs:}Azure oferece vários tipos de VMs para diferentes cargas de trabalho:
\begin{itemize}[left=0pt]

\item \textbf{Série B (Burstable)}: Projetada para cargas de trabalho intermitentes e de baixo custo.
\item \textbf{Série D (General Purpose)}: Equilibrada em CPU e memória, ideal para aplicativos de negócios.
\item \textbf{Série E (Memory Optimized)}: VMs com alta capacidade de memória, úteis para bancos de dados e cargas de trabalho de alto uso de memória.
\item \textbf{Série N (GPU)}: VMs com GPUs para processamento gráfico intenso, como renderização e aprendizado de máquina.
  
\end{itemize}

\subsection*{Formas de precificação}

\textbf{Pay-As-You-Go}: Pagamento por uso sem necessidade de compromissos de longo prazo. Ideal para ambientes de desenvolvimento e teste.
\textbf{Instâncias Reservadas:} Permite um compromisso de 1 ou 3 anos, resultando em até 72\% de desconto.

\subsection*{Azure SQL Database}

O Azure SQL Database é um serviço de banco de dados relacional gerenciado com compatibilidade SQL Server. Ele oferece diversas opções para escalabilidade e com suporte para \textit{HA} e  recuperação de desastres.

\subsection*{Forma de precificação}

\begin{itemize}

  \item Database Transation Unit - Ideal para cargas de trabalho com pequeno núbmero de transações.
  \item Baseado em em núbmero de Núcleos - Permite escalar o número de vCpu de forma independente.
  
\end{itemize}

\break

\section*{GCP}

\subsection*{Google Compute Engine (VMs)}

O Google Compute Engine fornece máquinas virtuais personalizáveis para hospedagem de aplicações com serviços de alta flexibilidade e escalabilidade.

\subsection*{Tipos de máquinas virtuais}
\begin{itemize}
  \item{E2} - Máquinas de usuo geral com equilíbrio entre custo e desempenho. 
  \item{N2} - Oferecem mais capacidade de CPU e memória recomendadas para cargas de trabalho de negócios mais exigentes.
  \item{M2} - Máquinas com alta quantidade de memória, adequada para grandes bancos de dados. 
  \item{A2} - Máquinas equipadas com GPUs para cargas envolvendo inteligência artificial.

\end{itemize}

\subsection*{Forma de precificação}

\begin{itemize}
  \item Commited use discounts - Reduz o cousto ao se comprometer com o uso durante um tempo determinado.
  \item Preemptible - Para cargas de trabalho tolerantes a falhas o custo é reduzido.
  \end{itemize}

\subsection*{Google Cloud Storage}
O Google Cloud Storage é oum serviço de armazenamento de objetos que oferece escalabilidade, segurança e alta disponibilidade para armazenar grandes volumes de dados, ele é ideal para backup, arquivamento e \textit{CDN}.
\subsection*{Classes de armazenamento}
O \textit{GCS} oferece quatro classes de armazenamento, com diferentes preços e níveis de disponibilidade de acesso aos dados:
\begin{itemize}
  \item\textbf{Standard} - Alta disponibilidade, ideal para acessar dados frequentemente.
  \item\textbf{Nearline} - Otimizado para dados que são acessados menos de uma vez por mês. 
  \item\textbf{Coldline} - Adequado para dados que não são acessados frequentemente, custando menos por GB.
  \item\textbf{Archive} - Para dados de longo prazo e raramente acessados, com custo reduzido.
\end{itemize}
\break

\section*{Alibaba Cloud}



\end{document}
